\chapter{Simulation Architecture for Semantic Galaxies}
\label{ch:simulation-architecture}

This chapter develops a complete computational framework for simulating
the RSVP--Polyxan dynamical system introduced in
Chapters~\ref{ch:embedding-flow}--\ref{ch:global-galaxy-maps}.
We treat the model as a \emph{physical analogue}: 
an interacting field--hypergraph--sheaf system whose evolution is governed
by geometric partial differential equations (PDEs), 
gradient flows on embeddings, and cohomological obstructions.

The goal of this chapter is to formalize a scalable simulation engine
that implements:

\begin{enumerate}
  \item the \emph{embedding flows}
        $\dot\iota = -\delta\mathcal{F}/\delta\iota$,
  \item the \emph{field evolution PDEs} for $(\Phi,\mathbf{v},S)$,
  \item the \emph{sheaf cohomology computations} that detect global semantic inconsistencies,
  \item and the \emph{feedback mechanisms} by which cohomology
        influences and stabilizes embedding dynamics.
\end{enumerate}

The resulting architecture is modular, parallelizable, and numerically stable.


% ============================================================
\section{Mathematical Structure of the Simulation}
% ============================================================

A simulation state at time $t$ is a tuple
\[
  \mathcal{X}(t)
  :=
  \Big(
    \iota_t : \mathcal{G}\to\mathcal{M},\;
    \Phi_t,\mathbf{v}_t,S_t,\;
    \Gamma_t,\;
    \mathfrak{o}(\iota_t)
  \Big),
\]
where:

\begin{itemize}
  \item $\mathcal{G}$ is the Polyxan hypergraph (finite or countable),
  \item $\iota_t$ is the embedding into $\mathcal{M}$,
  \item $(\Phi_t,\mathbf{v}_t,S_t)$ are RSVP fields on the discretized manifold,
  \item $\Gamma_t$ is a local galaxy map (partial or global),
  \item $\mathfrak{o}(\iota_t)$ is the obstruction class in sheaf cohomology.
\end{itemize}

The update of $\mathcal{X}(t)$ is governed by three coupled dynamical laws:

\begin{align*}
  \dot{\iota}_t &= - \frac{\delta \mathcal{F}}{\delta \iota}, \\
  \partial_t (\Phi,\mathbf{v},S)
    &= \mathcal{D}(\Phi,\mathbf{v},S) + \mathcal{J}(\iota), \\
  \mathfrak{o}(\iota_t)
    &= \mathrm{cohomology}(\iota_t, \mathcal{F},\mathscr{S}).
\end{align*}

Each term is formalized below.


% ============================================================
\section{Discretization of the Manifold and Hypergraph}
% ============================================================

\subsection{Manifold Discretization}

We discretize $\mathcal{M}$ using one of:

\begin{itemize}
  \item a simplicial mesh,
  \item a structured finite-difference grid,
  \item or a point cloud with neighborhood graph (for large scales).
\end{itemize}

Each mesh node carries:

\[
  \Big(\Phi(n),\mathbf{v}(n),S(n)\Big)
  \qquad n\in\mathcal{M}_h,
\]
and geometric coordinates $x(n)\in\mathbb{R}^d$.

Finite differences, finite elements, or spectral methods provide approximations
to:

\[
  \nabla,\quad \nabla^2,\quad \nabla\cdot\mathbf{v},\quad \mathrm{Ric}(g).
\]


\subsection{Hypergraph Representation}

The Polyxan hypergraph is encoded as:

\[
  \mathcal{G} = (V,E,w,m,\rho^\*),
\]
where:

\begin{itemize}
  \item $V$ is the set of nodes (atoms),
  \item $E$ is the set of hyperedges (relations),
  \item $w$ is semantic mass,
  \item $m$ is curvature mass (from Chapter~\ref{ch:typed-hypergraphs}),
  \item $\rho^\*$ is target density.
\end{itemize}

Local neighborhoods $U\subseteq\mathcal{G}$ form a cover for Čech cohomology.


\subsection{Embedding Storage}

An embedding is stored as a map:

\[
  \iota: V \to \mathcal{M}_h,
\]
where each hypergraph node maps to a manifold node or barycentric coordinates.

Embedding evolution requires:

\[
  \dot\iota(v) \in T_{\iota(v)}\mathcal{M},
\]
approximated using the discrete tangent space of the mesh.


% ============================================================
\section{RSVP Field PDEs}
% ============================================================

Field evolution is governed by the RSVP equations (Ch.~\ref{ch:rsvp-fields})
in discretized form:

\begin{align}
  \partial_t \Phi &= \nabla_h^2 \Phi 
                     + \gamma \nabla_h\cdot\mathbf{v} 
                     + J_{\Phi}^{(H)}
                     + \text{dissipation}, \\[0.2em]
  \partial_t S &= \beta \nabla_h \cdot (S^{-1}\nabla_h S)
                + J_S^{(H)}, \\[0.2em]
  \partial_t \mathbf{v}
    &= -\nabla_h \Phi
       + J_{\mathbf{v}}^{(H)}.
\end{align}

The hypergraph-derived sources $J^{(H)}$ are computed from:

\[
  J_\Phi^{(H)}(x(n))
  =
  \sum_{v\in V}
    m(v)\, \delta_{n,\iota(v)}.
\]

A standard explicit time-marching scheme is:

\[
  \Psi_{t+\Delta t}(n)
    = \Psi_{t}(n)
    + \Delta t\, \mathcal{D}(\Psi_t)(n)
    + \Delta t\,\mathcal{J}(\iota_t)(n),
\]
where $\Psi=(\Phi,\mathbf{v},S)$.


% ============================================================
\section{Embedding Gradient Flow}
% ============================================================

Recall the embedding energy:

\[
  \mathcal{F}
    = \mathcal{E}_{\mathrm{RSVP}}
    + \mathcal{E}_{\mathrm{curv}}
    + \mathcal{E}_{\mathrm{dens}}
    + \mathcal{E}_{\mathrm{sheaf}}.
\]

The discrete gradient flow is:

\[
  \iota_{t+\Delta t}(v)
  =
  \iota_t(v)
  - \Delta t\,
    \left(
      \frac{\delta\mathcal{E}_{\mathrm{RSVP}}}{\delta\iota(v)}
      +
      \frac{\delta\mathcal{E}_{\mathrm{curv}}}{\delta\iota(v)}
      +
      \frac{\delta\mathcal{E}_{\mathrm{dens}}}{\delta\iota(v)}
      +
      \frac{\delta\mathcal{E}_{\mathrm{sheaf}}}{\delta\iota(v)}
    \right).
\]

In mesh form:

\begin{align*}
  \frac{\delta\mathcal{E}_{\mathrm{RSVP}}}{\delta\iota(v)}
    &= w_v \nabla_h\mathcal{L}_{\mathrm{RSVP}}(\iota(v)), \\[0.3em]
  \frac{\delta\mathcal{E}_{\mathrm{curv}}}{\delta\iota(v)}
    &= 2\alpha_v H\,\partial H/\partial\iota(v), \\[0.3em]
  \frac{\delta\mathcal{E}_{\mathrm{dens}}}{\delta\iota(v)}
    &= 2\beta_v(\rho_\iota(v)-\rho^\*(v))\partial\rho_\iota/\partial\iota(v).
\end{align*}

The sheaf term is handled in the next section.


% ============================================================
\section{Cohomology, Galaxy Maps, and Feedback}
% ============================================================

\subsection{Local Sheaf Evaluation}

For each neighborhood $U_i$ in a cover:

\begin{itemize}
  \item restrict Polyxan sheaf $\mathcal{F}(U_i)$,
  \item restrict RSVP sheaf $\mathscr{S}(\iota(U_i))$,
  \item compute mismatch via candidate galaxy map $\Gamma_{U_i}$.
\end{itemize}

\subsection{Čech Complex Construction}

Construct cochains:

\[
  c^k_{i_0\cdots i_k}
  \in \mathscr{S}\big(\iota(U_{i_0}\cap\cdots\cap U_{i_k})\big),
\]
assembled into the Čech complex.

Compute obstruction class:

\[
  \mathfrak{o}(\iota)
  =
  [\delta c],
\]
where $\delta$ is the Čech coboundary.

\subsection{Feedback to Embeddings}

Define sheaf inconsistency energy:

\[
  \mathcal{E}_{\mathrm{sheaf}}
  =
  \sum_{U_i}
    \gamma_{U_i}\,
    \|c^1_{i,j}\|^2
  + \sum_{U_i,U_j,U_k}
    \gamma'_{i,j,k}\,
    \|c^2_{i,j,k}\|^2.
\]

Variation yields:

\[
  \frac{\delta\mathcal{E}_{\mathrm{sheaf}}}{\delta\iota(v)}
  =
  \sum_{U_i\ni v}
    \gamma_{U_i}\,
    \frac{\partial c^{1}_{i,j}}{\partial \iota(v)}
  +
  \sum_{U_i\ni v}
    \gamma'_{i,j,k}\,
    \frac{\partial c^{2}_{i,j,k}}{\partial \iota(v)}.
\]

Embedding flow is driven toward configurations minimizing cohomological defects.


% ============================================================
\section{Numerical Stability and Parallelization}
% ============================================================

\subsection{Stability}

Gradient flows require:

\[
  \Delta t < \frac{2}{\lambda_{\max}},
\]
where $\lambda_{\max}$ is the maximum eigenvalue of the Hessian of $\mathcal{F}$.

Field PDEs require analogous CFL conditions.

\subsection{Parallelization}

All major components parallelize:

\begin{itemize}
  \item RSVP PDEs over mesh nodes,
  \item hypergraph embedding updates over nodes $v\in V$,
  \item Čech cohomology over neighborhoods $U_i$.
\end{itemize}

GPU acceleration is effective for:

\begin{itemize}
  \item field Laplacians,
  \item convolution-like smoothing operators,
  \item sparse matrix operations for cohomology.
\end{itemize}


% ============================================================
\section{Interpretation 3 (Semantic) Layer}
% ============================================================

Although this chapter treats the simulation as a physical-analogue system,
the *same architecture* implements cognitive or semantic evolution:

\begin{itemize}
  \item RSVP fields become semantic potentials,
  \item Polyxan hypergraph nodes represent cognitive atoms,
  \item embeddings represent semantic placement in meaning-space,
  \item obstruction classes represent conceptual incoherence,
  \item gradient flows represent conceptual refinement.
\end{itemize}

Thus Interpretation~1 and Interpretation~3 share identical mathematics.


% ============================================================
\section{Simulation Loop Summary}
% ============================================================

At each time step:

\begin{enumerate}
\item Compute $\delta\mathcal{F}/\delta \iota$ for all $v\in V$.
\item Update embeddings $\iota \gets \iota - \Delta t\,\delta\mathcal{F}/\delta\iota$.
\item Update RSVP fields via discretized PDEs.
\item Compute Čech cochains and obstruction class.
\item Update $\mathcal{E}_{\mathrm{sheaf}}$ and feed back to embeddings.
\item Optionally visualize fields, embeddings, and obstruction maps.
\end{enumerate}

This defines the engine for semantic galaxy morphogenesis.


% ============================================================
\section{Summary}
% ============================================================

This chapter provided a complete simulation architecture for the RSVP--Polyxan
system, including:

\begin{itemize}
  \item discretization of manifold, fields, and hypergraphs,
  \item embedding gradient flows,
  \item RSVP PDE solvers,
  \item sheaf cohomology and obstruction feedback,
  \item numerical stability and parallelization,
  \item and a unified simulation loop.
\end{itemize}

Chapter~\ref{ch:nbody} will implement these ideas in an N-body
scheme and analyze emergent structure formation in semantic galaxies.
